\documentclass[12pt, letterpaper]{article}

\usepackage[utf8]{inputenc}
\title{Week 4}
\author{Raadhikka Gupta}
\date{\today}

\begin{document}
\maketitle

\section{Retrieval-Augmented Generation}
\subsection{Introduction}
\begin{itemize}
    \item RAG is an architecture for optimizing the performance of LLMs by connecting them to external knowledge bases; help LLMs deliver more relevant responses at a higher quality.
    \item Prompt $\rightarrow$ LLM $\rightarrow$ User Query $\rightarrow$ Information $\rightarrow$ LLM $\rightarrow$ User
    \item Use cases for RAG:
        \begin{itemize}
            \item \textbf{Question Answering}: can retrieve the latest and most relevant information from sources before generating answer.
            \item \textbf{Research Augmentation}: research accuracy is high as it combines historical data with recent trends to inform research findings.
            \item \textbf{Content Generation}: improves factual accuracy and overall credibility of content by including credible citations.
            \item \textbf{Domain-specific and proprietary data integration}: can retrieve information from internal cases without embedding it with training data and shared to public domain.
        \end{itemize}
\end{itemize}

\subsection{Benefits}
\begin{itemize}
    \item AI hallucination is a phenomenon in which an LLM detects patterns or objects that do not exist or are not visible to humans, leading to inaccurate or meaningless outputs.
    \item To confirm that organizations receive current, complete, and accurate responses is to retrain or fine-tune the model with updated data; but it is costly and impractical, so RAG is a solution to keep expanding its knowledge without retaining constantly.
    \item Without retraining, businesses can scale AI applications while reducing costs and resource demands.
    \item Allows to integrate with online APIs, social media feeds, and search engines to get information from external sources and stay up-to-date with market trends, news, and customer sentiment.
\end{itemize}

\subsection{RAG Workflow}
\begin{itemize}
    \item It consists of steps that combine information extracted from knowledge sources to produce more informed, precise, and relevant responses.
    \item \textbf{Prompt Creation}: submits a query that requires a response acting like a trigger that prompts RAG to get user input, determine the intent and scope of the information, and retrieve the relevant information.
    \item \textbf{Information Search}: searches external knowledge sources (research papers, books, company databases, APIs, web articles) which allows the response to the user be beyond the training dataset and most recent.
    \item \textbf{Information Retrieval}: filters and extracts the most relevant and credible information, making sure that only the high-quality data aligned with the user's query is selected.
    \item \textbf{Prompt Augmentation}: enhances the original query by integrating the retrieved information by providing the model with additional context to generate a more accurate and detailed answer aligned with the user's intent.
    \item \textbf{Response Generation}: uses enhanced information to generate a clear, accurate, and actionable response which is given to the user.
\end{itemize}

\end{document}